\subsection{Подраздел 1}

\td{Далее приводятся примеры для различных типов блоков.}

\ldots В качестве иллюстрации докажем следующие известные соотношения де~Моргана (см.~также рис.~\ref{fig:правила_де_моргана}) \ldots

\defi{defi:перестановки}
{Перестановки из $n$ элементов}{
    Всевозможные кортежи длины $n$, составленные из всех элементов данного $n$-элементного множества при условии, что каждый элемент входит в кортеж ровно один раз.
    Число всех возможных различных перестановок из $n$ элементов обозначается $P_n$.
}[
    \item Для множества $X = \set{1, 2, 3}$ все возможные перестановки~--- это кортежи $\tup{1, 2, 3}$, $\tup{1, 3, 2}$, $\tup{2, 1, 3}$, $\tup{2, 3, 1}$, $\tup{3, 1, 2}$, $\tup{3, 2, 1}$.
        Как можно видеть, в данном случае $P_3 = 6$.
]

\begin{figure}[t!]
    \centering
    \begin{subfigure}[b]{.45\textwidth}
        \centering \includegraphics[scale=0.58]{l-01/demorgan_union.pdf}
    \end{subfigure}
    \hfill
    \begin{subfigure}[b]{.45\textwidth}
        \centering \includegraphics[scale=0.58]{l-01/demorgan_intersection.pdf}
    \end{subfigure}
    \caption{Иллюстрация для правил де~Моргана.}
    \label{fig:правила_де_моргана}
\end{figure}

\theo{theo:перестановки}
{Количество перестановок из $n$ элементов}{
    \begin{equation}\label{eq:перестановки}
        \forall n \in \N \qs
        P_n = n!.
    \end{equation}
}[
    На первое место можно поставить любой из $n$ элементов.
    На второе~--- любой из оставшихся $n - 1$.
    На третье~--- любой из оставшихся $n - 2$, и~т.д.
    На последнее, $n$-е место, остаётся ровно один элемент.
    Тогда по правилу произведения:
    \begin{equation*}
        P_n = n \mul (n-1) \mul (n-2) \mul \ldots \mul 1 = n!.\qedhere
    \end{equation*}
][
    В данном разделе мы полагаем, что $n \ge 1$, но для пустого множества $X = \emptyset$, то есть для случая $n = 0$, тоже могут быть формально рассмотрены перестановки.
    При этом, очевидно, существует только одна возможная перестановка~--- пустой кортеж $\tup{}$.
    Заметим, что по формуле~\eqref{eq:перестановки} мы получаем аналогичный результат: $P_0 = 0! = 1$.
]

\exam{exam:перестановки_пример}
{Рассаживание студентов}
{
    В группе $6$ студентов.
    Сколькими способами можно рассадить всех студентов группы на $6$ стульях, стоящих в ряд?
}[
    Заметим, что каждый вариант рассаживания~--- это перестановка всех $6$ студентов.
    По формуле~\eqref{eq:перестановки} из теоремы~\ref{theo:перестановки} для количества перестановок имеем:
    \begin{equation*}
        P_6 = 6! = 720.\qedhere
    \end{equation*}
]

\lemm{lemm:размещения_и_сочетания}
{Связь размещений и сочетаний}{
    \begin{equation*}
        \forall n \in \N \qs
        \forall k \in \Nz, k \le n \qs
        A_n^k = C_n^k \mul P_k.
    \end{equation*}
}[
    Каждому $k$-элементному подмножеству (сочетанию) соответствует ровно $k!$ способов упорядочить его элементы (перестановок выбранных элементов).
    Поэтому число размещений равно произведению числа сочетаний на число способов упорядочить элементы каждого сочетания.
]

\note{note:функция_двух_переменных}
{Функция двух переменных}
{
    Функция двух переменных из множеств $X_1$ и $X_2$ в множество $Y$ является обычной функцией, область определения которой оказывается подмножеством прямого произведения $X_1 \times X_2$.
    Поэтому запись $F \colon X_1 \times X_2 \to Y$ означает, что каждой упорядоченной паре $(x_1, x_2)$ сопоставляется единственный элемент $F(x_1, x_2) \in Y$.
    Например, сложение натуральных чисел можно рассматривать как отображение $\mathord{+} \colon \N \times \N \to \N$, ставящее произвольной паре натуральных чисел $(m, n)$ в соответствие число $m + n$.
}


\stat{stat:отрицание_общего_высказывания}
{Отрицание общего высказывания}{
    \begin{equation*}
        \lnot (\forall x \in X \qs P(x))
        \eq
        \exists x_0 \in X \qs \lnot P(x_0).
    \end{equation*}
}